
\documentclass{article}
\usepackage{colortbl}
\usepackage{makecell}
\usepackage{multirow}
\usepackage{supertabular}

\begin{document}

\newcounter{utterance}

\centering \large Interaction Transcript for game `dond', experiment `semi\_en', episode 14 with Qwen3{-}32B{-}t0.0.
\vspace{24pt}

{ \footnotesize  \setcounter{utterance}{1}
\setlength{\tabcolsep}{0pt}
\begin{supertabular}{c@{$\;$}|p{.15\linewidth}@{}p{.15\linewidth}p{.15\linewidth}p{.15\linewidth}p{.15\linewidth}p{.15\linewidth}}
    \# & $\;$A & \multicolumn{4}{c}{Game Master} & $\;\:$B\\
    \hline

    \theutterance \stepcounter{utterance}  
    & & \multicolumn{4}{p{0.6\linewidth}}{
        \cellcolor[rgb]{0.9,0.9,0.9}{
            \makecell[{{p{\linewidth}}}]{
                \texttt{\tiny{[P1$\langle$GM]}}
                \texttt{You are playing a negotiation game in which you have to agree on how to divide a set of items among you and another player.} \\
\\ 
\texttt{Rules:} \\
\texttt{(a) You and the other player are given a set of items. Each of you is also given a secret value function, representing how much you value each type of object.} \\
\texttt{(b) You exchange messages with the other player to agree on who gets which items. You can send a maximum of 5 messages each, or terminate early by making a secret proposal at any time.} \\
\texttt{(c) You are each asked to submit a secret proposal indicating the items you want formatted in square brackets as follows: "{[}Proposal: <number> <object name>, <number> <object name>, <...>{]}"} \\
\texttt{(d) If your proposals are complementary, i.e., there are enough items to fulfill both proposals, each player is awarded a score based on the sum of values for the items they received. Otherwise, both of you get zero points.} \\
\texttt{(e) Your goal is to maximize the score you receive.} \\
\\ 
\texttt{Let us start.} \\
\\ 
\texttt{The set of available items is:} \\
\\ 
\texttt{1 frame, 1 shoe, 3 hammers} \\
\\ 
\texttt{Your secret value function is:} \\
\\ 
\texttt{frame: 8, shoe: 2, hammer: 0} \\
\\ 
\texttt{IMPORTANT: Your messages, unless it is the secret proposal, are directly transmitted to the other player, so do not include any response to the rules or text announcing your message. To make a secret proposal, use the indicated format. Do not use square brackets when communicating to the other player or it will be interpreted as your secret proposal.} \\
            }
        }
    }
    & \\ \\

    \theutterance \stepcounter{utterance}  
    & \multicolumn{4}{p{0.6\linewidth}}{
        \cellcolor[rgb]{0.9,0.9,0.9}{
            \makecell[{{p{\linewidth}}}]{
                \texttt{\tiny{[P1$\rangle$GM]}}
                \texttt{<think>} \\
\texttt{Okay, let's see. I need to negotiate the division of items between me and another player. The items available are 1 frame, 1 shoe, and 3 hammers. My value function is frame: 8, shoe: 2, hammer: 0. So, I value the frame the most, then the shoe, and hammers are worthless to me. The other player's value function is secret, so I don't know what they want. } \\
\\ 
\texttt{First, I should figure out what the other player might value. Since I don't have their value function, I need to make some assumptions. Maybe they value hammers more than me? Or maybe they also want the frame or the shoe. Since my goal is to maximize my own score, I need to get the items that give me the highest value. For me, that's the frame. The shoe is next, but it's only worth 2, and hammers are 0. So ideally, I want the frame and maybe the shoe, and leave the hammers to the other player.} \\
\\ 
\texttt{But I have to make sure that the proposals are complementary. If I take the frame and the shoe, then the other player would have the 3 hammers. If their value function is such that they value hammers, then that's good. But if they need the frame or the shoe, then our proposals might conflict. However, since I can't know their values, I have to make a proposal that is likely to be complementary. } \\
\\ 
\texttt{Wait, the rules say that if the proposals are complementary, meaning there are enough items to fulfill both, then we get our scores. Otherwise, both get zero. So I need to make a proposal that doesn't overlap with what the other player might want. Since I don't know their values, maybe I should propose taking the frame and the shoe, and leave the hammers. If the other player wants the hammers, then our proposals are complementary. If they also want the frame or shoe, then maybe not. But since I can't know, I have to assume that the other player might not value the frame and shoe as much as I do. } \\
\\ 
\texttt{Alternatively, maybe the other player values the frame and shoe more. But since my value for the frame is 8, which is high, I should try to secure it. Let me think. If I propose taking the frame, and leave the shoe and hammers to the} \\
            }
        }
    }
    & & \\ \\

    \theutterance \stepcounter{utterance}  
    & & & \multicolumn{2}{p{0.3\linewidth}}{
        \cellcolor[rgb]{0.9,0.9,0.9}{
            \makecell[{{p{\linewidth}}}]{
                \texttt{\tiny{[GM$|$GM]}}
                \texttt{continue} \\
            }
        }
    }
    & & \\ \\

    \theutterance \stepcounter{utterance}  
    & & & \multicolumn{2}{p{0.3\linewidth}}{
        \cellcolor[rgb]{0.9,0.9,0.9}{
            \makecell[{{p{\linewidth}}}]{
                \texttt{\tiny{[GM$|$GM]}}
                \texttt{<think>} \\
\texttt{Okay, let's see. I need to negotiate the division of items between me and another player. The items available are 1 frame, 1 shoe, and 3 hammers. My value function is frame: 8, shoe: 2, hammer: 0. So, I value the frame the most, then the shoe, and hammers are worthless to me. The other player's value function is secret, so I don't know what they want. } \\
\\ 
\texttt{First, I should figure out what the other player might value. Since I don't have their value function, I need to make some assumptions. Maybe they value hammers more than me? Or maybe they also want the frame or the shoe. Since my goal is to maximize my own score, I need to get the items that give me the highest value. For me, that's the frame. The shoe is next, but it's only worth 2, and hammers are 0. So ideally, I want the frame and maybe the shoe, and leave the hammers to the other player.} \\
\\ 
\texttt{But I have to make sure that the proposals are complementary. If I take the frame and the shoe, then the other player would have the 3 hammers. If their value function is such that they value hammers, then that's good. But if they need the frame or the shoe, then our proposals might conflict. However, since I can't know their values, I have to make a proposal that is likely to be complementary. } \\
\\ 
\texttt{Wait, the rules say that if the proposals are complementary, meaning there are enough items to fulfill both, then we get our scores. Otherwise, both get zero. So I need to make a proposal that doesn't overlap with what the other player might want. Since I don't know their values, maybe I should propose taking the frame and the shoe, and leave the hammers. If the other player wants the hammers, then our proposals are complementary. If they also want the frame or shoe, then maybe not. But since I can't know, I have to assume that the other player might not value the frame and shoe as much as I do. } \\
\\ 
\texttt{Alternatively, maybe the other player values the frame and shoe more. But since my value for the frame is 8, which is high, I should try to secure it. Let me think. If I propose taking the frame, and leave the shoe and hammers to the} \\
            }
        }
    }
    & & \\ \\

    \theutterance \stepcounter{utterance}  
    & & \multicolumn{4}{p{0.6\linewidth}}{
        \cellcolor[rgb]{0.9,0.9,0.9}{
            \makecell[{{p{\linewidth}}}]{
                \texttt{\tiny{[GM$\rangle$P2]}}
                \texttt{You are playing a negotiation game in which you have to agree on how to divide a set of items among you and another player.} \\
\\ 
\texttt{Rules:} \\
\texttt{(a) You and the other player are given a set of items. Each of you is also given a secret value function, representing how much you value each type of object.} \\
\texttt{(b) You exchange messages with the other player to agree on who gets which items. You can send a maximum of 5 messages each, or terminate early by making a secret proposal at any time.} \\
\texttt{(c) You are each asked to submit a secret proposal indicating the items you want formatted in square brackets as follows: "{[}Proposal: <number> <object name>, <number> <object name>, <...>{]}"} \\
\texttt{(d) If your proposals are complementary, i.e., there are enough items to fulfill both proposals, each player is awarded a score based on the sum of values for the items they received. Otherwise, both of you get zero points.} \\
\texttt{(e) Your goal is to maximize the score you receive.} \\
\\ 
\texttt{Let us start.} \\
\\ 
\texttt{The set of available items is:} \\
\\ 
\texttt{1 frame, 1 shoe, 3 hammers} \\
\\ 
\texttt{Your secret value function is:} \\
\\ 
\texttt{frame: 6, shoe: 1, hammer: 1} \\
\\ 
\texttt{IMPORTANT: Your messages, unless it is the secret proposal, are directly transmitted to the other player, so do not include any response to the rules or text announcing your message. To make a secret proposal, use the indicated format. Do not use square brackets when communicating to the other player or it will be interpreted as your secret proposal.} \\
\\ 
\\ 
\texttt{<think>} \\
\texttt{Okay, let's see. I need to negotiate the division of items between me and another player. The items available are 1 frame, 1 shoe, and 3 hammers. My value function is frame: 8, shoe: 2, hammer: 0. So, I value the frame the most, then the shoe, and hammers are worthless to me. The other player's value function is secret, so I don't know what they want. } \\
\\ 
\texttt{First, I should figure out what the other player might value. Since I don't have their value function, I need to make some assumptions. Maybe they value hammers more than me? Or maybe they also want the frame or the shoe. Since my goal is to maximize my own score, I need to get the items that give me the highest value. For me, that's the frame. The shoe is next, but it's only worth 2, and hammers are 0. So ideally, I want the frame and maybe the shoe, and leave the hammers to the other player.} \\
\\ 
\texttt{But I have to make sure that the proposals are complementary. If I take the frame and the shoe, then the other player would have the 3 hammers. If their value function is such that they value hammers, then that's good. But if they need the frame or the shoe, then our proposals might conflict. However, since I can't know their values, I have to make a proposal that is likely to be complementary. } \\
\\ 
\texttt{Wait, the rules say that if the proposals are complementary, meaning there are enough items to fulfill both, then we get our scores. Otherwise, both get zero. So I need to make a proposal that doesn't overlap with what the other player might want. Since I don't know their values, maybe I should propose taking the frame and the shoe, and leave the hammers. If the other player wants the hammers, then our proposals are complementary. If they also want the frame or shoe, then maybe not. But since I can't know, I have to assume that the other player might not value the frame and shoe as much as I do. } \\
\\ 
\texttt{Alternatively, maybe the other player values the frame and shoe more. But since my value for the frame is 8, which is high, I should try to secure it. Let me think. If I propose taking the frame, and leave the shoe and hammers to the} \\
            }
        }
    }
    & \\ \\

    \theutterance \stepcounter{utterance}  
    & & & \multicolumn{4}{p{0.6\linewidth}}{
        \cellcolor[rgb]{0.9,0.9,0.9}{
            \makecell[{{p{\linewidth}}}]{
                \texttt{\tiny{[GM$\langle$P2]}}
                \texttt{<think>} \\
\texttt{Okay, let's break this down. The user is in a negotiation game where they need to divide items between themselves and another player. The items are 1 frame, 1 shoe, and 3 hammers. Their own value function is frame: 8, shoe: 2, hammer: 0. The other player's values are secret. The goal is to maximize their own score by making a secret proposal that is complementary to the other player's proposal.} \\
\\ 
\texttt{First, the user's highest value is the frame. Since they get 8 points for it, they should prioritize getting the frame. The shoe is next, worth 2 points, but hammers are worthless. So ideally, they want the frame and maybe the shoe, and leave the hammers to the other player. But they need to ensure that their proposal doesn't conflict with the other player's. If the other player also wants the frame or shoe, their proposals might not be complementary. However, since the other player's values are unknown, the user has to make an educated guess.} \\
\\ 
\texttt{The user might consider that the other player could value hammers more, especially since there are three of them. If the user takes the frame and shoe, the other player could take all the hammers. That way, the proposals are complementary because there's no overlap. The user's total would be 8 + 2 = 10, and if the other player values hammers, they might get points from them. But if the other player also wants the frame or shoe, this could cause a conflict. However, without knowing the other's values, the user has to assume that the other might not value the frame and shoe as much. } \\
\\ 
\texttt{Alternatively, maybe the other player wants the frame. If the user takes the frame, the other player can't take it. But if the other player's highest value is the frame, then their proposals would conflict. But since the user's value for the frame is so high, they might be willing to risk that. However, if the other player also wants the frame, both proposals would conflict, leading to zero points for both. That's a risk. } \\
\\ 
\texttt{Another angle: the user could propose taking just the frame, leaving the shoe and hammers. If the other player takes the shoe and hammers, that's complementary. The user gets 8 points, and the other gets whatever their values are. But the user's shoe is only worth} \\
            }
        }
    }
     \\ \\

    \theutterance \stepcounter{utterance}  
    & & & \multicolumn{2}{p{0.3\linewidth}}{
        \cellcolor[rgb]{0.9,0.9,0.9}{
            \makecell[{{p{\linewidth}}}]{
                \texttt{\tiny{[GM$|$GM]}}
                \texttt{continue} \\
            }
        }
    }
    & & \\ \\

    \theutterance \stepcounter{utterance}  
    & & & \multicolumn{2}{p{0.3\linewidth}}{
        \cellcolor[rgb]{0.9,0.9,0.9}{
            \makecell[{{p{\linewidth}}}]{
                \texttt{\tiny{[GM$|$GM]}}
                \texttt{<think>} \\
\texttt{Okay, let's break this down. The user is in a negotiation game where they need to divide items between themselves and another player. The items are 1 frame, 1 shoe, and 3 hammers. Their own value function is frame: 8, shoe: 2, hammer: 0. The other player's values are secret. The goal is to maximize their own score by making a secret proposal that is complementary to the other player's proposal.} \\
\\ 
\texttt{First, the user's highest value is the frame. Since they get 8 points for it, they should prioritize getting the frame. The shoe is next, worth 2 points, but hammers are worthless. So ideally, they want the frame and maybe the shoe, and leave the hammers to the other player. But they need to ensure that their proposal doesn't conflict with the other player's. If the other player also wants the frame or shoe, their proposals might not be complementary. However, since the other player's values are unknown, the user has to make an educated guess.} \\
\\ 
\texttt{The user might consider that the other player could value hammers more, especially since there are three of them. If the user takes the frame and shoe, the other player could take all the hammers. That way, the proposals are complementary because there's no overlap. The user's total would be 8 + 2 = 10, and if the other player values hammers, they might get points from them. But if the other player also wants the frame or shoe, this could cause a conflict. However, without knowing the other's values, the user has to assume that the other might not value the frame and shoe as much. } \\
\\ 
\texttt{Alternatively, maybe the other player wants the frame. If the user takes the frame, the other player can't take it. But if the other player's highest value is the frame, then their proposals would conflict. But since the user's value for the frame is so high, they might be willing to risk that. However, if the other player also wants the frame, both proposals would conflict, leading to zero points for both. That's a risk. } \\
\\ 
\texttt{Another angle: the user could propose taking just the frame, leaving the shoe and hammers. If the other player takes the shoe and hammers, that's complementary. The user gets 8 points, and the other gets whatever their values are. But the user's shoe is only worth} \\
            }
        }
    }
    & & \\ \\

    \theutterance \stepcounter{utterance}  
    & & \multicolumn{4}{p{0.6\linewidth}}{
        \cellcolor[rgb]{0.9,0.9,0.9}{
            \makecell[{{p{\linewidth}}}]{
                \texttt{\tiny{[P1$\langle$GM]}}
                \texttt{<think>} \\
\texttt{Okay, let's break this down. The user is in a negotiation game where they need to divide items between themselves and another player. The items are 1 frame, 1 shoe, and 3 hammers. Their own value function is frame: 8, shoe: 2, hammer: 0. The other player's values are secret. The goal is to maximize their own score by making a secret proposal that is complementary to the other player's proposal.} \\
\\ 
\texttt{First, the user's highest value is the frame. Since they get 8 points for it, they should prioritize getting the frame. The shoe is next, worth 2 points, but hammers are worthless. So ideally, they want the frame and maybe the shoe, and leave the hammers to the other player. But they need to ensure that their proposal doesn't conflict with the other player's. If the other player also wants the frame or shoe, their proposals might not be complementary. However, since the other player's values are unknown, the user has to make an educated guess.} \\
\\ 
\texttt{The user might consider that the other player could value hammers more, especially since there are three of them. If the user takes the frame and shoe, the other player could take all the hammers. That way, the proposals are complementary because there's no overlap. The user's total would be 8 + 2 = 10, and if the other player values hammers, they might get points from them. But if the other player also wants the frame or shoe, this could cause a conflict. However, without knowing the other's values, the user has to assume that the other might not value the frame and shoe as much. } \\
\\ 
\texttt{Alternatively, maybe the other player wants the frame. If the user takes the frame, the other player can't take it. But if the other player's highest value is the frame, then their proposals would conflict. But since the user's value for the frame is so high, they might be willing to risk that. However, if the other player also wants the frame, both proposals would conflict, leading to zero points for both. That's a risk. } \\
\\ 
\texttt{Another angle: the user could propose taking just the frame, leaving the shoe and hammers. If the other player takes the shoe and hammers, that's complementary. The user gets 8 points, and the other gets whatever their values are. But the user's shoe is only worth} \\
            }
        }
    }
    & \\ \\

    \theutterance \stepcounter{utterance}  
    & \multicolumn{4}{p{0.6\linewidth}}{
        \cellcolor[rgb]{0.9,0.9,0.9}{
            \makecell[{{p{\linewidth}}}]{
                \texttt{\tiny{[P1$\rangle$GM]}}
                \texttt{<think>} \\
\texttt{Okay, let's try to figure out the best way to approach this. So, the user has to divide the items: 1 frame, 1 shoe, 3 hammers. Their own values are frame:8, shoe:2, hammer:0. The other player's values are secret. The goal is to maximize their own score by making a secret proposal that's complementary to the other player's.} \\
\\ 
\texttt{First, the user's highest value is the frame. So they really want the frame. The shoe is next, but it's only worth 2, and hammers are 0. So ideally, they want the frame and maybe the shoe, and leave the hammers to the other player. But they need to make sure that their proposal doesn't conflict with the other player's. If the other player also wants the frame or shoe, their proposals might not be complementary. However, since the other player's values are unknown, the user has to make an educated guess.} \\
\\ 
\texttt{The user might consider that the other player could value hammers more, especially since there are three of them. If the user takes the frame and shoe, the other player could take all the hammers. That way, the proposals are complementary because there's no overlap. The user's total would be 8 + 2 = 10, and if the other player values hammers, they might get points from them. But if the other player also wants the frame or shoe, this could cause a conflict. However, without knowing the other's values, the user has to assume that the other might not value the frame and shoe as much. } \\
\\ 
\texttt{Alternatively, maybe the other player wants the frame. If the user takes the frame, the other player can't take it. But if the other player's highest value is the frame, then their proposals would conflict. But since the user's value for the frame is so high, they might be willing to risk that. However, if the other player also wants the frame, both proposals would conflict, leading to zero points for both. That's a risk. } \\
\\ 
\texttt{Another angle: the user could propose taking just the frame, leaving the shoe and hammers. If the other player takes the shoe and hammers, that's complementary. The user gets 8 points, and the other gets whatever their values are. But the user's shoe is only worth 2, so maybe the other player values it more. But the user} \\
            }
        }
    }
    & & \\ \\

    \theutterance \stepcounter{utterance}  
    & & & \multicolumn{2}{p{0.3\linewidth}}{
        \cellcolor[rgb]{0.9,0.9,0.9}{
            \makecell[{{p{\linewidth}}}]{
                \texttt{\tiny{[GM$|$GM]}}
                \texttt{continue} \\
            }
        }
    }
    & & \\ \\

    \theutterance \stepcounter{utterance}  
    & & & \multicolumn{2}{p{0.3\linewidth}}{
        \cellcolor[rgb]{0.9,0.9,0.9}{
            \makecell[{{p{\linewidth}}}]{
                \texttt{\tiny{[GM$|$GM]}}
                \texttt{<think>} \\
\texttt{Okay, let's try to figure out the best way to approach this. So, the user has to divide the items: 1 frame, 1 shoe, 3 hammers. Their own values are frame:8, shoe:2, hammer:0. The other player's values are secret. The goal is to maximize their own score by making a secret proposal that's complementary to the other player's.} \\
\\ 
\texttt{First, the user's highest value is the frame. So they really want the frame. The shoe is next, but it's only worth 2, and hammers are 0. So ideally, they want the frame and maybe the shoe, and leave the hammers to the other player. But they need to make sure that their proposal doesn't conflict with the other player's. If the other player also wants the frame or shoe, their proposals might not be complementary. However, since the other player's values are unknown, the user has to make an educated guess.} \\
\\ 
\texttt{The user might consider that the other player could value hammers more, especially since there are three of them. If the user takes the frame and shoe, the other player could take all the hammers. That way, the proposals are complementary because there's no overlap. The user's total would be 8 + 2 = 10, and if the other player values hammers, they might get points from them. But if the other player also wants the frame or shoe, this could cause a conflict. However, without knowing the other's values, the user has to assume that the other might not value the frame and shoe as much. } \\
\\ 
\texttt{Alternatively, maybe the other player wants the frame. If the user takes the frame, the other player can't take it. But if the other player's highest value is the frame, then their proposals would conflict. But since the user's value for the frame is so high, they might be willing to risk that. However, if the other player also wants the frame, both proposals would conflict, leading to zero points for both. That's a risk. } \\
\\ 
\texttt{Another angle: the user could propose taking just the frame, leaving the shoe and hammers. If the other player takes the shoe and hammers, that's complementary. The user gets 8 points, and the other gets whatever their values are. But the user's shoe is only worth 2, so maybe the other player values it more. But the user} \\
            }
        }
    }
    & & \\ \\

    \theutterance \stepcounter{utterance}  
    & & \multicolumn{4}{p{0.6\linewidth}}{
        \cellcolor[rgb]{0.9,0.9,0.9}{
            \makecell[{{p{\linewidth}}}]{
                \texttt{\tiny{[GM$\rangle$P2]}}
                \texttt{<think>} \\
\texttt{Okay, let's try to figure out the best way to approach this. So, the user has to divide the items: 1 frame, 1 shoe, 3 hammers. Their own values are frame:8, shoe:2, hammer:0. The other player's values are secret. The goal is to maximize their own score by making a secret proposal that's complementary to the other player's.} \\
\\ 
\texttt{First, the user's highest value is the frame. So they really want the frame. The shoe is next, but it's only worth 2, and hammers are 0. So ideally, they want the frame and maybe the shoe, and leave the hammers to the other player. But they need to make sure that their proposal doesn't conflict with the other player's. If the other player also wants the frame or shoe, their proposals might not be complementary. However, since the other player's values are unknown, the user has to make an educated guess.} \\
\\ 
\texttt{The user might consider that the other player could value hammers more, especially since there are three of them. If the user takes the frame and shoe, the other player could take all the hammers. That way, the proposals are complementary because there's no overlap. The user's total would be 8 + 2 = 10, and if the other player values hammers, they might get points from them. But if the other player also wants the frame or shoe, this could cause a conflict. However, without knowing the other's values, the user has to assume that the other might not value the frame and shoe as much. } \\
\\ 
\texttt{Alternatively, maybe the other player wants the frame. If the user takes the frame, the other player can't take it. But if the other player's highest value is the frame, then their proposals would conflict. But since the user's value for the frame is so high, they might be willing to risk that. However, if the other player also wants the frame, both proposals would conflict, leading to zero points for both. That's a risk. } \\
\\ 
\texttt{Another angle: the user could propose taking just the frame, leaving the shoe and hammers. If the other player takes the shoe and hammers, that's complementary. The user gets 8 points, and the other gets whatever their values are. But the user's shoe is only worth 2, so maybe the other player values it more. But the user} \\
            }
        }
    }
    & \\ \\

    \theutterance \stepcounter{utterance}  
    & & & \multicolumn{4}{p{0.6\linewidth}}{
        \cellcolor[rgb]{0.9,0.9,0.9}{
            \makecell[{{p{\linewidth}}}]{
                \texttt{\tiny{[GM$\langle$P2]}}
                \texttt{<think>} \\
\texttt{Okay, let's try to figure out the best way to approach this. The user has to divide the items: 1 frame, 1 shoe, 3 hammers. Their own values are frame:8, shoe:2, hammer:0. The other player's values are secret. The goal is to maximize their own score by making a secret proposal that's complementary to the other player's.} \\
\\ 
\texttt{First, the user's highest value is the frame. So they really want the frame. The shoe is next, but it's only worth 2, and hammers are 0. So ideally, they want the frame and maybe the shoe, and leave the hammers to the other player. But they need to make sure that their proposal doesn't conflict with the other player's. If the other player also wants the frame or shoe, their proposals might not be complementary. However, since the other player's values are unknown, the user has to make an educated guess.} \\
\\ 
\texttt{The user might consider that the other player could value hammers more, especially since there are three of them. If the user takes the frame and shoe, the other player could take all the hammers. That way, the proposals are complementary because there's no overlap. The user's total would be 8 + 2 = 10, and if the other player values hammers, they might get points from them. But if the other player also wants the frame or shoe, this could cause a conflict. However, without knowing the other's values, the user has to assume that the other might not value the frame and shoe as much. } \\
\\ 
\texttt{Alternatively, maybe the other player wants the frame. If the user takes the frame, the other player can't take it. But if the other player's highest value is the frame, then their proposals would conflict. But since the user's value for the frame is so high, they might be willing to risk that. However, if the other player also wants the frame, both proposals would conflict, leading to zero points for both. That's a risk. } \\
\\ 
\texttt{Another angle: the user could propose taking just the frame, leaving the shoe and hammers. If the other player takes the shoe and hammers, that's complementary. The user gets 8 points, and the other gets whatever their values are. But the user's shoe is only worth 2, so maybe the other player values it more. But the user doesn't} \\
            }
        }
    }
     \\ \\

    \theutterance \stepcounter{utterance}  
    & & & \multicolumn{2}{p{0.3\linewidth}}{
        \cellcolor[rgb]{0.9,0.9,0.9}{
            \makecell[{{p{\linewidth}}}]{
                \texttt{\tiny{[GM$|$GM]}}
                \texttt{continue} \\
            }
        }
    }
    & & \\ \\

    \theutterance \stepcounter{utterance}  
    & & & \multicolumn{2}{p{0.3\linewidth}}{
        \cellcolor[rgb]{0.9,0.9,0.9}{
            \makecell[{{p{\linewidth}}}]{
                \texttt{\tiny{[GM$|$GM]}}
                \texttt{<think>} \\
\texttt{Okay, let's try to figure out the best way to approach this. The user has to divide the items: 1 frame, 1 shoe, 3 hammers. Their own values are frame:8, shoe:2, hammer:0. The other player's values are secret. The goal is to maximize their own score by making a secret proposal that's complementary to the other player's.} \\
\\ 
\texttt{First, the user's highest value is the frame. So they really want the frame. The shoe is next, but it's only worth 2, and hammers are 0. So ideally, they want the frame and maybe the shoe, and leave the hammers to the other player. But they need to make sure that their proposal doesn't conflict with the other player's. If the other player also wants the frame or shoe, their proposals might not be complementary. However, since the other player's values are unknown, the user has to make an educated guess.} \\
\\ 
\texttt{The user might consider that the other player could value hammers more, especially since there are three of them. If the user takes the frame and shoe, the other player could take all the hammers. That way, the proposals are complementary because there's no overlap. The user's total would be 8 + 2 = 10, and if the other player values hammers, they might get points from them. But if the other player also wants the frame or shoe, this could cause a conflict. However, without knowing the other's values, the user has to assume that the other might not value the frame and shoe as much. } \\
\\ 
\texttt{Alternatively, maybe the other player wants the frame. If the user takes the frame, the other player can't take it. But if the other player's highest value is the frame, then their proposals would conflict. But since the user's value for the frame is so high, they might be willing to risk that. However, if the other player also wants the frame, both proposals would conflict, leading to zero points for both. That's a risk. } \\
\\ 
\texttt{Another angle: the user could propose taking just the frame, leaving the shoe and hammers. If the other player takes the shoe and hammers, that's complementary. The user gets 8 points, and the other gets whatever their values are. But the user's shoe is only worth 2, so maybe the other player values it more. But the user doesn't} \\
            }
        }
    }
    & & \\ \\

    \theutterance \stepcounter{utterance}  
    & & \multicolumn{4}{p{0.6\linewidth}}{
        \cellcolor[rgb]{0.9,0.9,0.9}{
            \makecell[{{p{\linewidth}}}]{
                \texttt{\tiny{[P1$\langle$GM]}}
                \texttt{<think>} \\
\texttt{Okay, let's try to figure out the best way to approach this. The user has to divide the items: 1 frame, 1 shoe, 3 hammers. Their own values are frame:8, shoe:2, hammer:0. The other player's values are secret. The goal is to maximize their own score by making a secret proposal that's complementary to the other player's.} \\
\\ 
\texttt{First, the user's highest value is the frame. So they really want the frame. The shoe is next, but it's only worth 2, and hammers are 0. So ideally, they want the frame and maybe the shoe, and leave the hammers to the other player. But they need to make sure that their proposal doesn't conflict with the other player's. If the other player also wants the frame or shoe, their proposals might not be complementary. However, since the other player's values are unknown, the user has to make an educated guess.} \\
\\ 
\texttt{The user might consider that the other player could value hammers more, especially since there are three of them. If the user takes the frame and shoe, the other player could take all the hammers. That way, the proposals are complementary because there's no overlap. The user's total would be 8 + 2 = 10, and if the other player values hammers, they might get points from them. But if the other player also wants the frame or shoe, this could cause a conflict. However, without knowing the other's values, the user has to assume that the other might not value the frame and shoe as much. } \\
\\ 
\texttt{Alternatively, maybe the other player wants the frame. If the user takes the frame, the other player can't take it. But if the other player's highest value is the frame, then their proposals would conflict. But since the user's value for the frame is so high, they might be willing to risk that. However, if the other player also wants the frame, both proposals would conflict, leading to zero points for both. That's a risk. } \\
\\ 
\texttt{Another angle: the user could propose taking just the frame, leaving the shoe and hammers. If the other player takes the shoe and hammers, that's complementary. The user gets 8 points, and the other gets whatever their values are. But the user's shoe is only worth 2, so maybe the other player values it more. But the user doesn't} \\
            }
        }
    }
    & \\ \\

    \theutterance \stepcounter{utterance}  
    & \multicolumn{4}{p{0.6\linewidth}}{
        \cellcolor[rgb]{0.9,0.9,0.9}{
            \makecell[{{p{\linewidth}}}]{
                \texttt{\tiny{[P1$\rangle$GM]}}
                \texttt{<think>} \\
\texttt{Okay, let's try to figure out the best way to approach this. The user has to divide the items: 1 frame, 1 shoe, 3 hammers. Their own values are frame:8, shoe:2, hammer:0. The other player's values are secret. The goal is to maximize their own score by making a secret proposal that's complementary to the other player's.} \\
\\ 
\texttt{First, the user's highest value is the frame. So they really want the frame. The shoe is next, but it's only worth 2, and hammers are 0. So ideally, they want the frame and maybe the shoe, and leave the hammers to the other player. But they need to make sure that their proposal doesn't conflict with the other player's. If the other player also wants the frame or shoe, their proposals might not be complementary. However, since the other player's values are unknown, the user has to make an educated guess.} \\
\\ 
\texttt{The user might consider that the other player could value hammers more, especially since there are three of them. If the user takes the frame and shoe, the other player could take all the hammers. That way, the proposals are complementary because there's no overlap. The user's total would be 8 + 2 = 10, and if the other player values hammers, they might get points from them. But if the other player also wants the frame or shoe, this could cause a conflict. However, without knowing the other's values, the user has to assume that the other might not value the frame and shoe as much. } \\
\\ 
\texttt{Alternatively, maybe the other player wants the frame. If the user takes the frame, the other player can't take it. But if the other player's highest value is the frame, then their proposals would conflict. But since the user's value for the frame is so high, they might be willing to risk that. However, if the other player also wants the frame, both proposals would conflict, leading to zero points for both. That's a risk. } \\
\\ 
\texttt{Another angle: the user could propose taking just the frame, leaving the shoe and hammers. If the other player takes the shoe and hammers, that's complementary. The user gets 8 points, and the other gets whatever their values are. But the user's shoe is only worth 2, so maybe the other player values it more. But the user doesn't} \\
            }
        }
    }
    & & \\ \\

    \theutterance \stepcounter{utterance}  
    & & & \multicolumn{2}{p{0.3\linewidth}}{
        \cellcolor[rgb]{0.9,0.9,0.9}{
            \makecell[{{p{\linewidth}}}]{
                \texttt{\tiny{[GM$|$GM]}}
                \texttt{continue} \\
            }
        }
    }
    & & \\ \\

    \theutterance \stepcounter{utterance}  
    & & & \multicolumn{2}{p{0.3\linewidth}}{
        \cellcolor[rgb]{0.9,0.9,0.9}{
            \makecell[{{p{\linewidth}}}]{
                \texttt{\tiny{[GM$|$GM]}}
                \texttt{<think>} \\
\texttt{Okay, let's try to figure out the best way to approach this. The user has to divide the items: 1 frame, 1 shoe, 3 hammers. Their own values are frame:8, shoe:2, hammer:0. The other player's values are secret. The goal is to maximize their own score by making a secret proposal that's complementary to the other player's.} \\
\\ 
\texttt{First, the user's highest value is the frame. So they really want the frame. The shoe is next, but it's only worth 2, and hammers are 0. So ideally, they want the frame and maybe the shoe, and leave the hammers to the other player. But they need to make sure that their proposal doesn't conflict with the other player's. If the other player also wants the frame or shoe, their proposals might not be complementary. However, since the other player's values are unknown, the user has to make an educated guess.} \\
\\ 
\texttt{The user might consider that the other player could value hammers more, especially since there are three of them. If the user takes the frame and shoe, the other player could take all the hammers. That way, the proposals are complementary because there's no overlap. The user's total would be 8 + 2 = 10, and if the other player values hammers, they might get points from them. But if the other player also wants the frame or shoe, this could cause a conflict. However, without knowing the other's values, the user has to assume that the other might not value the frame and shoe as much. } \\
\\ 
\texttt{Alternatively, maybe the other player wants the frame. If the user takes the frame, the other player can't take it. But if the other player's highest value is the frame, then their proposals would conflict. But since the user's value for the frame is so high, they might be willing to risk that. However, if the other player also wants the frame, both proposals would conflict, leading to zero points for both. That's a risk. } \\
\\ 
\texttt{Another angle: the user could propose taking just the frame, leaving the shoe and hammers. If the other player takes the shoe and hammers, that's complementary. The user gets 8 points, and the other gets whatever their values are. But the user's shoe is only worth 2, so maybe the other player values it more. But the user doesn't} \\
            }
        }
    }
    & & \\ \\

    \theutterance \stepcounter{utterance}  
    & & \multicolumn{4}{p{0.6\linewidth}}{
        \cellcolor[rgb]{0.9,0.9,0.9}{
            \makecell[{{p{\linewidth}}}]{
                \texttt{\tiny{[GM$\rangle$P2]}}
                \texttt{<think>} \\
\texttt{Okay, let's try to figure out the best way to approach this. The user has to divide the items: 1 frame, 1 shoe, 3 hammers. Their own values are frame:8, shoe:2, hammer:0. The other player's values are secret. The goal is to maximize their own score by making a secret proposal that's complementary to the other player's.} \\
\\ 
\texttt{First, the user's highest value is the frame. So they really want the frame. The shoe is next, but it's only worth 2, and hammers are 0. So ideally, they want the frame and maybe the shoe, and leave the hammers to the other player. But they need to make sure that their proposal doesn't conflict with the other player's. If the other player also wants the frame or shoe, their proposals might not be complementary. However, since the other player's values are unknown, the user has to make an educated guess.} \\
\\ 
\texttt{The user might consider that the other player could value hammers more, especially since there are three of them. If the user takes the frame and shoe, the other player could take all the hammers. That way, the proposals are complementary because there's no overlap. The user's total would be 8 + 2 = 10, and if the other player values hammers, they might get points from them. But if the other player also wants the frame or shoe, this could cause a conflict. However, without knowing the other's values, the user has to assume that the other might not value the frame and shoe as much. } \\
\\ 
\texttt{Alternatively, maybe the other player wants the frame. If the user takes the frame, the other player can't take it. But if the other player's highest value is the frame, then their proposals would conflict. But since the user's value for the frame is so high, they might be willing to risk that. However, if the other player also wants the frame, both proposals would conflict, leading to zero points for both. That's a risk. } \\
\\ 
\texttt{Another angle: the user could propose taking just the frame, leaving the shoe and hammers. If the other player takes the shoe and hammers, that's complementary. The user gets 8 points, and the other gets whatever their values are. But the user's shoe is only worth 2, so maybe the other player values it more. But the user doesn't} \\
            }
        }
    }
    & \\ \\

    \theutterance \stepcounter{utterance}  
    & & & \multicolumn{4}{p{0.6\linewidth}}{
        \cellcolor[rgb]{0.9,0.9,0.9}{
            \makecell[{{p{\linewidth}}}]{
                \texttt{\tiny{[GM$\langle$P2]}}
                \texttt{<think>} \\
\texttt{Okay, let's try to figure out the best way to approach this. The user has to divide the items: 1 frame, 1 shoe, 3 hammers. Their own values are frame:8, shoe:2, hammer:0. The other player's values are secret. The goal is to maximize their own score by making a secret proposal that's complementary to the other player's.} \\
\\ 
\texttt{First, the user's highest value is the frame. So they really want the frame. The shoe is next, but it's only worth 2, and hammers are 0. So ideally, they want the frame and maybe the shoe, and leave the hammers to the other player. But they need to make sure that their proposal doesn't conflict with the other player's. If the other player also wants the frame or shoe, their proposals might not be complementary. However, since the other player's values are unknown, the user has to make an educated guess.} \\
\\ 
\texttt{The user might consider that the other player could value hammers more, especially since there are three of them. If the user takes the frame and shoe, the other player could take all the hammers. That way, the proposals are complementary because there's no overlap. The user's total would be 8 + 2 = 10, and if the other player values hammers, they might get points from them. But if the other player also wants the frame or shoe, this could cause a conflict. However, without knowing the other's values, the user has to assume that the other might not value the frame and shoe as much. } \\
\\ 
\texttt{Alternatively, maybe the other player wants the frame. If the user takes the frame, the other player can't take it. But if the other player's highest value is the frame, then their proposals would conflict. But since the user's value for the frame is so high, they might be willing to risk that. However, if the other player also wants the frame, both proposals would conflict, leading to zero points for both. That's a risk. } \\
\\ 
\texttt{Another angle: the user could propose taking just the frame, leaving the shoe and hammers. If the other player takes the shoe and hammers, that's complementary. The user gets 8 points, and the other gets whatever their values are. But the user's shoe is only worth 2, so maybe the other player values it more. But the user doesn't} \\
            }
        }
    }
     \\ \\

    \theutterance \stepcounter{utterance}  
    & & & \multicolumn{2}{p{0.3\linewidth}}{
        \cellcolor[rgb]{0.9,0.9,0.9}{
            \makecell[{{p{\linewidth}}}]{
                \texttt{\tiny{[GM$|$GM]}}
                \texttt{continue} \\
            }
        }
    }
    & & \\ \\

    \theutterance \stepcounter{utterance}  
    & & & \multicolumn{2}{p{0.3\linewidth}}{
        \cellcolor[rgb]{0.9,0.9,0.9}{
            \makecell[{{p{\linewidth}}}]{
                \texttt{\tiny{[GM$|$GM]}}
                \texttt{<think>} \\
\texttt{Okay, let's try to figure out the best way to approach this. The user has to divide the items: 1 frame, 1 shoe, 3 hammers. Their own values are frame:8, shoe:2, hammer:0. The other player's values are secret. The goal is to maximize their own score by making a secret proposal that's complementary to the other player's.} \\
\\ 
\texttt{First, the user's highest value is the frame. So they really want the frame. The shoe is next, but it's only worth 2, and hammers are 0. So ideally, they want the frame and maybe the shoe, and leave the hammers to the other player. But they need to make sure that their proposal doesn't conflict with the other player's. If the other player also wants the frame or shoe, their proposals might not be complementary. However, since the other player's values are unknown, the user has to make an educated guess.} \\
\\ 
\texttt{The user might consider that the other player could value hammers more, especially since there are three of them. If the user takes the frame and shoe, the other player could take all the hammers. That way, the proposals are complementary because there's no overlap. The user's total would be 8 + 2 = 10, and if the other player values hammers, they might get points from them. But if the other player also wants the frame or shoe, this could cause a conflict. However, without knowing the other's values, the user has to assume that the other might not value the frame and shoe as much. } \\
\\ 
\texttt{Alternatively, maybe the other player wants the frame. If the user takes the frame, the other player can't take it. But if the other player's highest value is the frame, then their proposals would conflict. But since the user's value for the frame is so high, they might be willing to risk that. However, if the other player also wants the frame, both proposals would conflict, leading to zero points for both. That's a risk. } \\
\\ 
\texttt{Another angle: the user could propose taking just the frame, leaving the shoe and hammers. If the other player takes the shoe and hammers, that's complementary. The user gets 8 points, and the other gets whatever their values are. But the user's shoe is only worth 2, so maybe the other player values it more. But the user doesn't} \\
            }
        }
    }
    & & \\ \\

    \theutterance \stepcounter{utterance}  
    & & \multicolumn{4}{p{0.6\linewidth}}{
        \cellcolor[rgb]{0.9,0.9,0.9}{
            \makecell[{{p{\linewidth}}}]{
                \texttt{\tiny{[P1$\langle$GM]}}
                \texttt{<think>} \\
\texttt{Okay, let's try to figure out the best way to approach this. The user has to divide the items: 1 frame, 1 shoe, 3 hammers. Their own values are frame:8, shoe:2, hammer:0. The other player's values are secret. The goal is to maximize their own score by making a secret proposal that's complementary to the other player's.} \\
\\ 
\texttt{First, the user's highest value is the frame. So they really want the frame. The shoe is next, but it's only worth 2, and hammers are 0. So ideally, they want the frame and maybe the shoe, and leave the hammers to the other player. But they need to make sure that their proposal doesn't conflict with the other player's. If the other player also wants the frame or shoe, their proposals might not be complementary. However, since the other player's values are unknown, the user has to make an educated guess.} \\
\\ 
\texttt{The user might consider that the other player could value hammers more, especially since there are three of them. If the user takes the frame and shoe, the other player could take all the hammers. That way, the proposals are complementary because there's no overlap. The user's total would be 8 + 2 = 10, and if the other player values hammers, they might get points from them. But if the other player also wants the frame or shoe, this could cause a conflict. However, without knowing the other's values, the user has to assume that the other might not value the frame and shoe as much. } \\
\\ 
\texttt{Alternatively, maybe the other player wants the frame. If the user takes the frame, the other player can't take it. But if the other player's highest value is the frame, then their proposals would conflict. But since the user's value for the frame is so high, they might be willing to risk that. However, if the other player also wants the frame, both proposals would conflict, leading to zero points for both. That's a risk. } \\
\\ 
\texttt{Another angle: the user could propose taking just the frame, leaving the shoe and hammers. If the other player takes the shoe and hammers, that's complementary. The user gets 8 points, and the other gets whatever their values are. But the user's shoe is only worth 2, so maybe the other player values it more. But the user doesn't} \\
            }
        }
    }
    & \\ \\

    \theutterance \stepcounter{utterance}  
    & \multicolumn{4}{p{0.6\linewidth}}{
        \cellcolor[rgb]{0.9,0.9,0.9}{
            \makecell[{{p{\linewidth}}}]{
                \texttt{\tiny{[P1$\rangle$GM]}}
                \texttt{<think>} \\
\texttt{Okay, let's try to figure out the best way to approach this. The user has to divide the items: 1 frame, 1 shoe, 3 hammers. Their own values are frame:8, shoe:2, hammer:0. The other player's values are secret. The goal is to maximize their own score by making a secret proposal that's complementary to the other player's.} \\
\\ 
\texttt{First, the user's highest value is the frame. So they really want the frame. The shoe is next, but it's only worth 2, and hammers are 0. So ideally, they want the frame and maybe the shoe, and leave the hammers to the other player. But they need to make sure that their proposal doesn't conflict with the other player's. If the other player also wants the frame or shoe, their proposals might not be complementary. However, since the other player's values are unknown, the user has to make an educated guess.} \\
\\ 
\texttt{The user might consider that the other player could value hammers more, especially since there are three of them. If the user takes the frame and shoe, the other player could take all the hammers. That way, the proposals are complementary because there's no overlap. The user's total would be 8 + 2 = 10, and if the other player values hammers, they might get points from them. But if the other player also wants the frame or shoe, this could cause a conflict. However, without knowing the other's values, the user has to assume that the other might not value the frame and shoe as much. } \\
\\ 
\texttt{Alternatively, maybe the other player wants the frame. If the user takes the frame, the other player can't take it. But if the other player's highest value is the frame, then their proposals would conflict. But since the user's value for the frame is so high, they might be willing to risk that. However, if the other player also wants the frame, both proposals would conflict, leading to zero points for both. That's a risk. } \\
\\ 
\texttt{Another angle: the user could propose taking just the frame, leaving the shoe and hammers. If the other player takes the shoe and hammers, that's complementary. The user gets 8 points, and the other gets whatever their values are. But the user's shoe is only worth 2, so maybe the other player values it more. But the user doesn't} \\
            }
        }
    }
    & & \\ \\

    \theutterance \stepcounter{utterance}  
    & & & \multicolumn{2}{p{0.3\linewidth}}{
        \cellcolor[rgb]{0.9,0.9,0.9}{
            \makecell[{{p{\linewidth}}}]{
                \texttt{\tiny{[GM$|$GM]}}
                \texttt{continue} \\
            }
        }
    }
    & & \\ \\

    \theutterance \stepcounter{utterance}  
    & & & \multicolumn{2}{p{0.3\linewidth}}{
        \cellcolor[rgb]{0.9,0.9,0.9}{
            \makecell[{{p{\linewidth}}}]{
                \texttt{\tiny{[GM$|$GM]}}
                \texttt{<think>} \\
\texttt{Okay, let's try to figure out the best way to approach this. The user has to divide the items: 1 frame, 1 shoe, 3 hammers. Their own values are frame:8, shoe:2, hammer:0. The other player's values are secret. The goal is to maximize their own score by making a secret proposal that's complementary to the other player's.} \\
\\ 
\texttt{First, the user's highest value is the frame. So they really want the frame. The shoe is next, but it's only worth 2, and hammers are 0. So ideally, they want the frame and maybe the shoe, and leave the hammers to the other player. But they need to make sure that their proposal doesn't conflict with the other player's. If the other player also wants the frame or shoe, their proposals might not be complementary. However, since the other player's values are unknown, the user has to make an educated guess.} \\
\\ 
\texttt{The user might consider that the other player could value hammers more, especially since there are three of them. If the user takes the frame and shoe, the other player could take all the hammers. That way, the proposals are complementary because there's no overlap. The user's total would be 8 + 2 = 10, and if the other player values hammers, they might get points from them. But if the other player also wants the frame or shoe, this could cause a conflict. However, without knowing the other's values, the user has to assume that the other might not value the frame and shoe as much. } \\
\\ 
\texttt{Alternatively, maybe the other player wants the frame. If the user takes the frame, the other player can't take it. But if the other player's highest value is the frame, then their proposals would conflict. But since the user's value for the frame is so high, they might be willing to risk that. However, if the other player also wants the frame, both proposals would conflict, leading to zero points for both. That's a risk. } \\
\\ 
\texttt{Another angle: the user could propose taking just the frame, leaving the shoe and hammers. If the other player takes the shoe and hammers, that's complementary. The user gets 8 points, and the other gets whatever their values are. But the user's shoe is only worth 2, so maybe the other player values it more. But the user doesn't} \\
            }
        }
    }
    & & \\ \\

    \theutterance \stepcounter{utterance}  
    & & \multicolumn{4}{p{0.6\linewidth}}{
        \cellcolor[rgb]{0.9,0.9,0.9}{
            \makecell[{{p{\linewidth}}}]{
                \texttt{\tiny{[GM$\rangle$P2]}}
                \texttt{<think>} \\
\texttt{Okay, let's try to figure out the best way to approach this. The user has to divide the items: 1 frame, 1 shoe, 3 hammers. Their own values are frame:8, shoe:2, hammer:0. The other player's values are secret. The goal is to maximize their own score by making a secret proposal that's complementary to the other player's.} \\
\\ 
\texttt{First, the user's highest value is the frame. So they really want the frame. The shoe is next, but it's only worth 2, and hammers are 0. So ideally, they want the frame and maybe the shoe, and leave the hammers to the other player. But they need to make sure that their proposal doesn't conflict with the other player's. If the other player also wants the frame or shoe, their proposals might not be complementary. However, since the other player's values are unknown, the user has to make an educated guess.} \\
\\ 
\texttt{The user might consider that the other player could value hammers more, especially since there are three of them. If the user takes the frame and shoe, the other player could take all the hammers. That way, the proposals are complementary because there's no overlap. The user's total would be 8 + 2 = 10, and if the other player values hammers, they might get points from them. But if the other player also wants the frame or shoe, this could cause a conflict. However, without knowing the other's values, the user has to assume that the other might not value the frame and shoe as much. } \\
\\ 
\texttt{Alternatively, maybe the other player wants the frame. If the user takes the frame, the other player can't take it. But if the other player's highest value is the frame, then their proposals would conflict. But since the user's value for the frame is so high, they might be willing to risk that. However, if the other player also wants the frame, both proposals would conflict, leading to zero points for both. That's a risk. } \\
\\ 
\texttt{Another angle: the user could propose taking just the frame, leaving the shoe and hammers. If the other player takes the shoe and hammers, that's complementary. The user gets 8 points, and the other gets whatever their values are. But the user's shoe is only worth 2, so maybe the other player values it more. But the user doesn't} \\
            }
        }
    }
    & \\ \\

    \theutterance \stepcounter{utterance}  
    & & & \multicolumn{4}{p{0.6\linewidth}}{
        \cellcolor[rgb]{0.9,0.9,0.9}{
            \makecell[{{p{\linewidth}}}]{
                \texttt{\tiny{[GM$\langle$P2]}}
                \texttt{<think>} \\
\texttt{Okay, let's try to figure out the best way to approach this. The user has to divide the items: 1 frame, 1 shoe, 3 hammers. Their own values are frame:8, shoe:2, hammer:0. The other player's values are secret. The goal is to maximize their own score by making a secret proposal that's complementary to the other player's.} \\
\\ 
\texttt{First, the user's highest value is the frame. So they really want the frame. The shoe is next, but it's only worth 2, and hammers are 0. So ideally, they want the frame and maybe the shoe, and leave the hammers to the other player. But they need to make sure that their proposal doesn't conflict with the other player's. If the other player also wants the frame or shoe, their proposals might not be complementary. However, since the other player's values are unknown, the user has to make an educated guess.} \\
\\ 
\texttt{The user might consider that the other player could value hammers more, especially since there are three of them. If the user takes the frame and shoe, the other player could take all the hammers. That way, the proposals are complementary because there's no overlap. The user's total would be 8 + 2 = 10, and if the other player values hammers, they might get points from them. But if the other player also wants the frame or shoe, this could cause a conflict. However, without knowing the other's values, the user has to assume that the other might not value the frame and shoe as much. } \\
\\ 
\texttt{Alternatively, maybe the other player wants the frame. If the user takes the frame, the other player can't take it. But if the other player's highest value is the frame, then their proposals would conflict. But since the user's value for the frame is so high, they might be willing to risk that. However, if the other player also wants the frame, both proposals would conflict, leading to zero points for both. That's a risk. } \\
\\ 
\texttt{Another angle: the user could propose taking just the frame, leaving the shoe and hammers. If the other player takes the shoe and hammers, that's complementary. The user gets 8 points, and the other gets whatever their values are. But the user's shoe is only worth 2, so maybe the other player values it more. But the user doesn't} \\
            }
        }
    }
     \\ \\

    \theutterance \stepcounter{utterance}  
    & & & \multicolumn{2}{p{0.3\linewidth}}{
        \cellcolor[rgb]{0.9,0.9,0.9}{
            \makecell[{{p{\linewidth}}}]{
                \texttt{\tiny{[GM$|$GM]}}
                \texttt{continue} \\
            }
        }
    }
    & & \\ \\

    \theutterance \stepcounter{utterance}  
    & & & \multicolumn{2}{p{0.3\linewidth}}{
        \cellcolor[rgb]{0.9,0.9,0.9}{
            \makecell[{{p{\linewidth}}}]{
                \texttt{\tiny{[GM$|$GM]}}
                \texttt{<think>} \\
\texttt{Okay, let's try to figure out the best way to approach this. The user has to divide the items: 1 frame, 1 shoe, 3 hammers. Their own values are frame:8, shoe:2, hammer:0. The other player's values are secret. The goal is to maximize their own score by making a secret proposal that's complementary to the other player's.} \\
\\ 
\texttt{First, the user's highest value is the frame. So they really want the frame. The shoe is next, but it's only worth 2, and hammers are 0. So ideally, they want the frame and maybe the shoe, and leave the hammers to the other player. But they need to make sure that their proposal doesn't conflict with the other player's. If the other player also wants the frame or shoe, their proposals might not be complementary. However, since the other player's values are unknown, the user has to make an educated guess.} \\
\\ 
\texttt{The user might consider that the other player could value hammers more, especially since there are three of them. If the user takes the frame and shoe, the other player could take all the hammers. That way, the proposals are complementary because there's no overlap. The user's total would be 8 + 2 = 10, and if the other player values hammers, they might get points from them. But if the other player also wants the frame or shoe, this could cause a conflict. However, without knowing the other's values, the user has to assume that the other might not value the frame and shoe as much. } \\
\\ 
\texttt{Alternatively, maybe the other player wants the frame. If the user takes the frame, the other player can't take it. But if the other player's highest value is the frame, then their proposals would conflict. But since the user's value for the frame is so high, they might be willing to risk that. However, if the other player also wants the frame, both proposals would conflict, leading to zero points for both. That's a risk. } \\
\\ 
\texttt{Another angle: the user could propose taking just the frame, leaving the shoe and hammers. If the other player takes the shoe and hammers, that's complementary. The user gets 8 points, and the other gets whatever their values are. But the user's shoe is only worth 2, so maybe the other player values it more. But the user doesn't} \\
            }
        }
    }
    & & \\ \\

    \theutterance \stepcounter{utterance}  
    & & \multicolumn{4}{p{0.6\linewidth}}{
        \cellcolor[rgb]{0.9,0.9,0.9}{
            \makecell[{{p{\linewidth}}}]{
                \texttt{\tiny{[P1$\langle$GM]}}
                \texttt{<think>} \\
\texttt{Okay, let's try to figure out the best way to approach this. The user has to divide the items: 1 frame, 1 shoe, 3 hammers. Their own values are frame:8, shoe:2, hammer:0. The other player's values are secret. The goal is to maximize their own score by making a secret proposal that's complementary to the other player's.} \\
\\ 
\texttt{First, the user's highest value is the frame. So they really want the frame. The shoe is next, but it's only worth 2, and hammers are 0. So ideally, they want the frame and maybe the shoe, and leave the hammers to the other player. But they need to make sure that their proposal doesn't conflict with the other player's. If the other player also wants the frame or shoe, their proposals might not be complementary. However, since the other player's values are unknown, the user has to make an educated guess.} \\
\\ 
\texttt{The user might consider that the other player could value hammers more, especially since there are three of them. If the user takes the frame and shoe, the other player could take all the hammers. That way, the proposals are complementary because there's no overlap. The user's total would be 8 + 2 = 10, and if the other player values hammers, they might get points from them. But if the other player also wants the frame or shoe, this could cause a conflict. However, without knowing the other's values, the user has to assume that the other might not value the frame and shoe as much. } \\
\\ 
\texttt{Alternatively, maybe the other player wants the frame. If the user takes the frame, the other player can't take it. But if the other player's highest value is the frame, then their proposals would conflict. But since the user's value for the frame is so high, they might be willing to risk that. However, if the other player also wants the frame, both proposals would conflict, leading to zero points for both. That's a risk. } \\
\\ 
\texttt{Another angle: the user could propose taking just the frame, leaving the shoe and hammers. If the other player takes the shoe and hammers, that's complementary. The user gets 8 points, and the other gets whatever their values are. But the user's shoe is only worth 2, so maybe the other player values it more. But the user doesn't} \\
            }
        }
    }
    & \\ \\

    \theutterance \stepcounter{utterance}  
    & \multicolumn{4}{p{0.6\linewidth}}{
        \cellcolor[rgb]{0.9,0.9,0.9}{
            \makecell[{{p{\linewidth}}}]{
                \texttt{\tiny{[P1$\rangle$GM]}}
                \texttt{<think>} \\
\texttt{Okay, let's try to figure out the best way to approach this. The user has to divide the items: 1 frame, 1 shoe, 3 hammers. Their own values are frame:8, shoe:2, hammer:0. The other player's values are secret. The goal is to maximize their own score by making a secret proposal that's complementary to the other player's.} \\
\\ 
\texttt{First, the user's highest value is the frame. So they really want the frame. The shoe is next, but it's only worth 2, and hammers are 0. So ideally, they want the frame and maybe the shoe, and leave the hammers to the other player. But they need to make sure that their proposal doesn't conflict with the other player's. If the other player also wants the frame or shoe, their proposals might not be complementary. However, since the other player's values are unknown, the user has to make an educated guess.} \\
\\ 
\texttt{The user might consider that the other player could value hammers more, especially since there are three of them. If the user takes the frame and shoe, the other player could take all the hammers. That way, the proposals are complementary because there's no overlap. The user's total would be 8 + 2 = 10, and if the other player values hammers, they might get points from them. But if the other player also wants the frame or shoe, this could cause a conflict. However, without knowing the other's values, the user has to assume that the other might not value the frame and shoe as much. } \\
\\ 
\texttt{Alternatively, maybe the other player wants the frame. If the user takes the frame, the other player can't take it. But if the other player's highest value is the frame, then their proposals would conflict. But since the user's value for the frame is so high, they might be willing to risk that. However, if the other player also wants the frame, both proposals would conflict, leading to zero points for both. That's a risk. } \\
\\ 
\texttt{Another angle: the user could propose taking just the frame, leaving the shoe and hammers. If the other player takes the shoe and hammers, that's complementary. The user gets 8 points, and the other gets whatever their values are. But the user's shoe is only worth 2, so maybe the other player values it more. But the user doesn't} \\
            }
        }
    }
    & & \\ \\

    \theutterance \stepcounter{utterance}  
    & & & \multicolumn{2}{p{0.3\linewidth}}{
        \cellcolor[rgb]{0.9,0.9,0.9}{
            \makecell[{{p{\linewidth}}}]{
                \texttt{\tiny{[GM$|$GM]}}
                \texttt{continue} \\
            }
        }
    }
    & & \\ \\

    \theutterance \stepcounter{utterance}  
    & & & \multicolumn{2}{p{0.3\linewidth}}{
        \cellcolor[rgb]{0.9,0.9,0.9}{
            \makecell[{{p{\linewidth}}}]{
                \texttt{\tiny{[GM$|$GM]}}
                \texttt{<think>} \\
\texttt{Okay, let's try to figure out the best way to approach this. The user has to divide the items: 1 frame, 1 shoe, 3 hammers. Their own values are frame:8, shoe:2, hammer:0. The other player's values are secret. The goal is to maximize their own score by making a secret proposal that's complementary to the other player's.} \\
\\ 
\texttt{First, the user's highest value is the frame. So they really want the frame. The shoe is next, but it's only worth 2, and hammers are 0. So ideally, they want the frame and maybe the shoe, and leave the hammers to the other player. But they need to make sure that their proposal doesn't conflict with the other player's. If the other player also wants the frame or shoe, their proposals might not be complementary. However, since the other player's values are unknown, the user has to make an educated guess.} \\
\\ 
\texttt{The user might consider that the other player could value hammers more, especially since there are three of them. If the user takes the frame and shoe, the other player could take all the hammers. That way, the proposals are complementary because there's no overlap. The user's total would be 8 + 2 = 10, and if the other player values hammers, they might get points from them. But if the other player also wants the frame or shoe, this could cause a conflict. However, without knowing the other's values, the user has to assume that the other might not value the frame and shoe as much. } \\
\\ 
\texttt{Alternatively, maybe the other player wants the frame. If the user takes the frame, the other player can't take it. But if the other player's highest value is the frame, then their proposals would conflict. But since the user's value for the frame is so high, they might be willing to risk that. However, if the other player also wants the frame, both proposals would conflict, leading to zero points for both. That's a risk. } \\
\\ 
\texttt{Another angle: the user could propose taking just the frame, leaving the shoe and hammers. If the other player takes the shoe and hammers, that's complementary. The user gets 8 points, and the other gets whatever their values are. But the user's shoe is only worth 2, so maybe the other player values it more. But the user doesn't} \\
            }
        }
    }
    & & \\ \\

    \theutterance \stepcounter{utterance}  
    & & \multicolumn{4}{p{0.6\linewidth}}{
        \cellcolor[rgb]{0.9,0.9,0.9}{
            \makecell[{{p{\linewidth}}}]{
                \texttt{\tiny{[GM$\rangle$P2]}}
                \texttt{<think>} \\
\texttt{Okay, let's try to figure out the best way to approach this. The user has to divide the items: 1 frame, 1 shoe, 3 hammers. Their own values are frame:8, shoe:2, hammer:0. The other player's values are secret. The goal is to maximize their own score by making a secret proposal that's complementary to the other player's.} \\
\\ 
\texttt{First, the user's highest value is the frame. So they really want the frame. The shoe is next, but it's only worth 2, and hammers are 0. So ideally, they want the frame and maybe the shoe, and leave the hammers to the other player. But they need to make sure that their proposal doesn't conflict with the other player's. If the other player also wants the frame or shoe, their proposals might not be complementary. However, since the other player's values are unknown, the user has to make an educated guess.} \\
\\ 
\texttt{The user might consider that the other player could value hammers more, especially since there are three of them. If the user takes the frame and shoe, the other player could take all the hammers. That way, the proposals are complementary because there's no overlap. The user's total would be 8 + 2 = 10, and if the other player values hammers, they might get points from them. But if the other player also wants the frame or shoe, this could cause a conflict. However, without knowing the other's values, the user has to assume that the other might not value the frame and shoe as much. } \\
\\ 
\texttt{Alternatively, maybe the other player wants the frame. If the user takes the frame, the other player can't take it. But if the other player's highest value is the frame, then their proposals would conflict. But since the user's value for the frame is so high, they might be willing to risk that. However, if the other player also wants the frame, both proposals would conflict, leading to zero points for both. That's a risk. } \\
\\ 
\texttt{Another angle: the user could propose taking just the frame, leaving the shoe and hammers. If the other player takes the shoe and hammers, that's complementary. The user gets 8 points, and the other gets whatever their values are. But the user's shoe is only worth 2, so maybe the other player values it more. But the user doesn't} \\
            }
        }
    }
    & \\ \\

    \theutterance \stepcounter{utterance}  
    & & & \multicolumn{4}{p{0.6\linewidth}}{
        \cellcolor[rgb]{0.9,0.9,0.9}{
            \makecell[{{p{\linewidth}}}]{
                \texttt{\tiny{[GM$\langle$P2]}}
                \texttt{<think>} \\
\texttt{Okay, let's try to figure out the best way to approach this. The user has to divide the items: 1 frame, 1 shoe, 3 hammers. Their own values are frame:8, shoe:2, hammer:0. The other player's values are secret. The goal is to maximize their own score by making a secret proposal that's complementary to the other player's.} \\
\\ 
\texttt{First, the user's highest value is the frame. So they really want the frame. The shoe is next, but it's only worth 2, and hammers are 0. So ideally, they want the frame and maybe the shoe, and leave the hammers to the other player. But they need to make sure that their proposal doesn't conflict with the other player's. If the other player also wants the frame or shoe, their proposals might not be complementary. However, since the other player's values are unknown, the user has to make an educated guess.} \\
\\ 
\texttt{The user might consider that the other player could value hammers more, especially since there are three of them. If the user takes the frame and shoe, the other player could take all the hammers. That way, the proposals are complementary because there's no overlap. The user's total would be 8 + 2 = 10, and if the other player values hammers, they might get points from them. But if the other player also wants the frame or shoe, this could cause a conflict. However, without knowing the other's values, the user has to assume that the other might not value the frame and shoe as much. } \\
\\ 
\texttt{Alternatively, maybe the other player wants the frame. If the user takes the frame, the other player can't take it. But if the other player's highest value is the frame, then their proposals would conflict. But since the user's value for the frame is so high, they might be willing to risk that. However, if the other player also wants the frame, both proposals would conflict, leading to zero points for both. That's a risk. } \\
\\ 
\texttt{Another angle: the user could propose taking just the frame, leaving the shoe and hammers. If the other player takes the shoe and hammers, that's complementary. The user gets 8 points, and the other gets whatever their values are. But the user's shoe is only worth 2, so maybe the other player values it more. But the user doesn't} \\
            }
        }
    }
     \\ \\

    \theutterance \stepcounter{utterance}  
    & & & \multicolumn{2}{p{0.3\linewidth}}{
        \cellcolor[rgb]{0.9,0.9,0.9}{
            \makecell[{{p{\linewidth}}}]{
                \texttt{\tiny{[GM$|$GM]}}
                \texttt{continue} \\
            }
        }
    }
    & & \\ \\

    \theutterance \stepcounter{utterance}  
    & & & \multicolumn{2}{p{0.3\linewidth}}{
        \cellcolor[rgb]{0.9,0.9,0.9}{
            \makecell[{{p{\linewidth}}}]{
                \texttt{\tiny{[GM$|$GM]}}
                \texttt{<think>} \\
\texttt{Okay, let's try to figure out the best way to approach this. The user has to divide the items: 1 frame, 1 shoe, 3 hammers. Their own values are frame:8, shoe:2, hammer:0. The other player's values are secret. The goal is to maximize their own score by making a secret proposal that's complementary to the other player's.} \\
\\ 
\texttt{First, the user's highest value is the frame. So they really want the frame. The shoe is next, but it's only worth 2, and hammers are 0. So ideally, they want the frame and maybe the shoe, and leave the hammers to the other player. But they need to make sure that their proposal doesn't conflict with the other player's. If the other player also wants the frame or shoe, their proposals might not be complementary. However, since the other player's values are unknown, the user has to make an educated guess.} \\
\\ 
\texttt{The user might consider that the other player could value hammers more, especially since there are three of them. If the user takes the frame and shoe, the other player could take all the hammers. That way, the proposals are complementary because there's no overlap. The user's total would be 8 + 2 = 10, and if the other player values hammers, they might get points from them. But if the other player also wants the frame or shoe, this could cause a conflict. However, without knowing the other's values, the user has to assume that the other might not value the frame and shoe as much. } \\
\\ 
\texttt{Alternatively, maybe the other player wants the frame. If the user takes the frame, the other player can't take it. But if the other player's highest value is the frame, then their proposals would conflict. But since the user's value for the frame is so high, they might be willing to risk that. However, if the other player also wants the frame, both proposals would conflict, leading to zero points for both. That's a risk. } \\
\\ 
\texttt{Another angle: the user could propose taking just the frame, leaving the shoe and hammers. If the other player takes the shoe and hammers, that's complementary. The user gets 8 points, and the other gets whatever their values are. But the user's shoe is only worth 2, so maybe the other player values it more. But the user doesn't} \\
            }
        }
    }
    & & \\ \\

    \theutterance \stepcounter{utterance}  
    & & \multicolumn{4}{p{0.6\linewidth}}{
        \cellcolor[rgb]{0.9,0.9,0.9}{
            \makecell[{{p{\linewidth}}}]{
                \texttt{\tiny{[P1$\langle$GM]}}
                \texttt{<think>} \\
\texttt{Okay, let's try to figure out the best way to approach this. The user has to divide the items: 1 frame, 1 shoe, 3 hammers. Their own values are frame:8, shoe:2, hammer:0. The other player's values are secret. The goal is to maximize their own score by making a secret proposal that's complementary to the other player's.} \\
\\ 
\texttt{First, the user's highest value is the frame. So they really want the frame. The shoe is next, but it's only worth 2, and hammers are 0. So ideally, they want the frame and maybe the shoe, and leave the hammers to the other player. But they need to make sure that their proposal doesn't conflict with the other player's. If the other player also wants the frame or shoe, their proposals might not be complementary. However, since the other player's values are unknown, the user has to make an educated guess.} \\
\\ 
\texttt{The user might consider that the other player could value hammers more, especially since there are three of them. If the user takes the frame and shoe, the other player could take all the hammers. That way, the proposals are complementary because there's no overlap. The user's total would be 8 + 2 = 10, and if the other player values hammers, they might get points from them. But if the other player also wants the frame or shoe, this could cause a conflict. However, without knowing the other's values, the user has to assume that the other might not value the frame and shoe as much. } \\
\\ 
\texttt{Alternatively, maybe the other player wants the frame. If the user takes the frame, the other player can't take it. But if the other player's highest value is the frame, then their proposals would conflict. But since the user's value for the frame is so high, they might be willing to risk that. However, if the other player also wants the frame, both proposals would conflict, leading to zero points for both. That's a risk. } \\
\\ 
\texttt{Another angle: the user could propose taking just the frame, leaving the shoe and hammers. If the other player takes the shoe and hammers, that's complementary. The user gets 8 points, and the other gets whatever their values are. But the user's shoe is only worth 2, so maybe the other player values it more. But the user doesn't} \\
\\ 
\\ 
\texttt{The time is up. It is now your turn to submit a secret proposal.} \\
            }
        }
    }
    & \\ \\

    \theutterance \stepcounter{utterance}  
    & \multicolumn{4}{p{0.6\linewidth}}{
        \cellcolor[rgb]{0.9,0.9,0.9}{
            \makecell[{{p{\linewidth}}}]{
                \texttt{\tiny{[P1$\rangle$GM]}}
                \texttt{<think>} \\
\texttt{Okay, let's try to figure out the best way to approach this. The user has to divide the items: 1 frame, 1 shoe, 3 hammers. Their own values are frame:8, shoe:2, hammer:0. The other player's values are secret. The goal is to maximize their own score by making a secret proposal that's complementary to the other player's.} \\
\\ 
\texttt{First, the user's highest value is the frame. So they really want the frame. The shoe is next, but it's only worth 2, and hammers are 0. So ideally, they want the frame and maybe the shoe, and leave the hammers to the other player. But they need to make sure that their proposal doesn't conflict with the other player's. If the other player also wants the frame or shoe, their proposals might not be complementary. However, since the other player's values are unknown, the user has to make an educated guess.} \\
\\ 
\texttt{The user might consider that the other player could value hammers more, especially since there are three of them. If the user takes the frame and shoe, the other player could take all the hammers. That way, the proposals are complementary because there's no overlap. The user's total would be 8 + 2 = 10, and if the other player values hammers, they might get points from them. But if the other player also wants the frame or shoe, this could cause a conflict. However, without knowing the other's values, the user has to assume that the other might not value the frame and shoe as much. } \\
\\ 
\texttt{Alternatively, maybe the other player wants the frame. If the user takes the frame, the other player can't take it. But if the other player's highest value is the frame, then their proposals would conflict. But since the user's value for the frame is so high, they might be willing to risk that. However, if the other player also wants the frame, both proposals would conflict, leading to zero points for both. That's a risk. } \\
\\ 
\texttt{Another angle: the user could propose taking just the frame, leaving the shoe and hammers. If the other player takes the shoe and hammers, that's complementary. The user gets 8 points, and the other gets whatever their values are. But the user's shoe is only worth 2, so maybe the other player values it more. But the user doesn't} \\
            }
        }
    }
    & & \\ \\

    \theutterance \stepcounter{utterance}  
    & & & \multicolumn{2}{p{0.3\linewidth}}{
        \cellcolor[rgb]{0.9,0.9,0.9}{
            \makecell[{{p{\linewidth}}}]{
                \texttt{\tiny{[GM$|$GM]}}
                \texttt{continue} \\
            }
        }
    }
    & & \\ \\

    \theutterance \stepcounter{utterance}  
    & & & \multicolumn{2}{p{0.3\linewidth}}{
        \cellcolor[rgb]{0.9,0.9,0.9}{
            \makecell[{{p{\linewidth}}}]{
                \texttt{\tiny{[GM$|$GM]}}
                \texttt{player was instructed to make a proposal} \\
            }
        }
    }
    & & \\ \\

\end{supertabular}
}

\end{document}
